\section {The QHFT functor}\label{TENSOR} Consider a QHFT bordism
$\Bb =(\tilde{Y},\tilde{L}_\Ff,\rho,\alpha_\pm)$. For every odd
integer $N\geq 1$, we associate a finite dimensional complex
linear space $V(\alpha_\pm)$ to $\alpha_{\pm}$, and a linear map
$\Hh_N(\Bb) : V(\alpha_-) \to V(\alpha_+)$ to $\Bb$, well-defined
up to sign and multiplication by $N$th roots of unity. This
defines a (moderately projective) functor $\Hh_N: {\rm QHB}
\rightarrow {\rm Vect}$, where Vect is the tensor category of
complex linear spaces. The construction immediately implies that
$\Hh_N$ is a modular functor, in the sense of \cite{T}, III.1.2.

\subsection {From QHFT bordisms to QHG-triangulated pseudomanifolds}
\label{Z(B)}
First we associate to $\Bb$ a pseudomanifold $Z(\Bb)$.
Fill each tunnel boundary component of $\tilde{Y}$
with a solid tube, thus recovering a copy of the manifold $Y$. The
cores of the solid tubes make a parallel unframed copy $\lambda''$
of $L$. We define $Z(\Bb)$ as the result of collapsing to one
point each component of $\lambda''$. In other words, we glue to
each tunnel component of $\tilde{Y}$ the {\it oriented topological
cusp} $\hat{\bf C} = B \times [0,+\infty]/(B \times \{\infty\})$
with base equal to either $B=S^1\times [-1,1]$ or $B=S^1\times
S^1$.
\smallskip

Next we describe a procedure to convert $Z(\Bb)$ to a
distinguished QHG-triangulated pseudomanifold. We refer to the
notions introduced in Subsection \ref{ToGeoSpec}.

\smallskip

We say that a branched triangulation $(T,b)$ of $B$ as above is {\it
  admissible} if $B \cap \partial_\pm \tilde{Y}$ and the tangles
  $\lambda$, $\lambda'$ are covered by the $1$-skeleton. We denote
  $(\hat{T},\hat{b})$ the branched triangulation of $\hat{\bf C}$,
  where $\hat{T}$ is the cone over $T$ from the non manifold point,
  say $\infty$, and $\hat{b}$ extends $b$ so that $\infty$ is a pit
  for every branched tetrahedron of $\hat{T}$.  Assume we are given an
  idealizable $PB^+(2,\mc)$-valued cocycle $z$ on an admissible
  triangulation of $B$. The idealization of $z$ determines for each
  $2$-simplex of $B$ a face of an ideal hyperbolic tetrahedron with
  further vertex at $\infty$ (see Figure \ref{tcusp}, where opposite
  vertical triangles are identified).  Since the fundamental group of
  $B$ is Abelian, the resulting family of oriented ideal hyperbolic
  tetrahedra actually makes an $\Ii$-triangulation
  $(\hat{T},\hat{b},w)$, which we call an {\it $\Ii$-cusp}. By
  conjugating if necessary, we see that $\Ii$-cusps make sense also
  when $z$ takes values more generally in $PSL(2,\mc)$. We get
  flattenings similarly as in Lemma \ref{canonicalflat}: at a corner
  of a $2$-simplex formed by edges $e_l$ and $e_r$ we put the
  difference of the logarithms of the vectors in $\mc$ associated by
  the idealization to $e_l$ and $e_r$, respectively.
\begin{figure}[ht]
\begin{center}
\includegraphics[width=4cm]{tcusp.eps}
\caption{\label{tcusp} An $\Ii$-cusp.}
\end{center}
\end{figure}

\begin{defi}\label{Dtriang}{\rm A {\it $\Dd$-triangulation} of
$\Bb =(\tilde{Y},\tilde{L}_\Ff,\rho,\alpha_\pm)$ consists of a
$4$-uple $\Kk =(K,\bar{H},b,z)$ where:

\smallskip

(a) $(K,b)$ is a branched triangulation of $\tilde{Y}$ extending
that on $\partial_\pm \tilde{Y}$, and inducing an admissible cusp
base triangulation at each tunnel component of $\tilde{L}_{\Ff}$.

\smallskip

(b) The $1$-dimensional subcomplex $\bar{H}= H\cup H'$ of $K$ is
ambiently isotopic to the tangle $\bar{\lambda}=\lambda \cup
\lambda'$, and $\bar{H}$ contains
all the vertices of $K$.

\smallskip

(c) $z$ is an idealizable $PSL(2,\C)$-valued $1$-cocycle on $(K,b)$
such that:

(i) the conjugacy class of $PSL(2,\mc)$-representations of $\pi_1(\tilde{Y})$
associated to $z$ coincides with $\rho$;

(ii) the ($+$)-exponential $\Ii$-parameters given by the
restriction of $z$ to $\partial_{\pm} \tilde{Y}$ coincide with
that of the objects $\alpha_- \cup \alpha_+$ (see Definition \ref{log-I}).
\smallskip

(d) the restriction of $z$ to each vertical tunnel component of
$\tilde{L}_{\Ff}$ takes values in the Borel subgroup $PB^+(2,\mc)$
of $PSL(2,\mc)$.}
\end{defi}
For any $\Dd$-triangulation $\Kk =(K,\bar{H},b,z)$ of $\Bb$ we get
a distinguished $\Ii$-triangulation $\Kk_\Ii=(K,\bar{H},b,w)$ of
$(Z(\Bb), \bar{\lambda})$ by gluing the idealization of $\Kk$ with
the $\Ii$-cusp given by the cocycle at each tunnel component. Note
that $\bar{H}$ contains all the vertices of $K_I$ that are
manifold points.
\begin{defi}\label{flat/charge-def} {\rm We say that $\mathcal{T}(\mathcal{B})
=(\Kk_\Ii,f,c) = (K,\bar{H},b,w,f,c)$ is a {\it distinguished
  flat/charged $\Ii$-triangulation} of $(Z(\Bb),\bar{\lambda})$ if it
satisfies Definition \ref{DISFLATCHARGE}, and at every boundary edge
of $Z(\Bb)$ the total (classical or quantum) log-branch of
$(\Kk_\Ii,f,c)$ of (\ref{totw}) coincides with the (classical or
quantum) Log-$\Ii$-parameter of the boundary object $\alpha_- \cup
\alpha_+$. }
\end{defi}
Recall the cohomological weights from section \ref{ToGeoSpec}.  These
notions still make sense for distinguished flat/charged
$\Ii$-triangulations, where the homology of $\partial T_0$ is replaced
with that of the tunnel components $\tilde{L}_\Ff$. Since we have
weights $(h_f,k_f)$ for log-branches and $(h_c,k_c)$ for charges at
the same time, we will denote them $(h,k)=((h_f,h_c),(k_f,k_c))$. We
can also define, in the very same way, boundary weights $k_f\in
H^1(\partial \tilde{Y}; \mc)$, but these are completely encoded by
$\alpha_\pm$. We have:
\begin{teo} \label{teoTOP} For every bordism
$\Bb=(\tilde{Y},\tilde{L}_\Ff,\rho,\alpha_\pm)$ and every $(h,k)
\in H^1(\tilde{Y};\Z/2\Z) \times H^1(\tilde{L}_\Ff; \mc)$
satisfying the properties (\ref{weightcondition}), there are
distinguished flat/charged $\Ii$-triangulations $\Tt(\Bb)$ of
$Z(\Bb)$ with weight $(h,k)$, and any two are QHG-isomorphic.
\end{teo}

\noindent {\it Proof.} The existence of distinguished
$\Ii$-triangulations $\Kk_\Ii$ of $Z(\Bb)$ follows from a tedious
but straightforward generalization of Theorem 4.13 in \cite{BB1}.
Global flattenings and integral charges with arbitrary weight
exist on the double $D\Kk_\Ii$ of $\Kk_\Ii$ by the results
recalled with Theorem \ref{theoweight}.

Consider the ($+$)-Log-$\Ii$-parameters $\{W^\pm(z)(e)\}_e$ at
$\alpha_\pm$. They are in one-one correspondence with the interior
edges of the corresponding $e$-triangulations, which is less than the
cardinality of the families $m$ and $n$ used to defined flattenings
and charges in Definition \ref{log-I}. Hence any family of
determinations of the logarithms of the $W^\pm(z)(e)$ is a system of
($+$)-Log-$\Ii$-parameters. Also, Lemma \ref{canonicalflat} implies
that any system of ($+$)-Log-$\Ii$-parameters at $\alpha_\pm$ extends
to a distinguished flat/charge $\Ii$-triangulation of the pseudo
manifold obtained from the trivial cylinders over
$\partial_{\pm}\tilde{Y}$ by collapsing to a point each annulus of
$\partial( \partial_{\pm}\tilde{Y}) \times [-1,1]$.

This means that any QHFT surface bounds a QHFT bordism, and that
for the bordism $\Bb$ there are flat/charges on $D\Kk_\Ii$ whose
restriction to $\Kk_\Ii$ induce the Log-$\Ii$-parameters of
$\alpha_\pm$. Hence we get global flat/charges as in Definition
\ref{flat/charge-def}. In fact the affine spaces of flat/charges
on $D\Kk_\Ii$ project onto that on $\Kk_\Ii$ compatible with
$\alpha_\pm$ (see the end of section \ref{QHG}). Then, the
Mayer-Vietoris exact sequence in cohomology for the triad
$(DK_\Ii,\Kk_\Ii,-\Kk_\Ii)$ shows that $(h,k)$ is induced by some
weight on $D\Kk_\Ii$. As we can choose the latter arbitrarily,
this concludes the proof of the first claim.

The second is harder, but follows strictly from the arguments in
the proof of Theorem 6.8 (2) in \cite{BB2}. The only new
ingredient is the presence of $\Ii$-cusps, which mimic the ends of
cusped manifolds treated in that paper. \hfill$\Box$

An alternative characterization of classical/quantum log-branches
of distinguished flat/charged $\Ii$-triangulations $\Tt(\Bb)$
follows from Remark \ref{univfc}.
\subsection{Amplitudes}\label{QHFTTENS} Fix an odd positive integer $N$.
Write $V=\C^N$, with the canonical basis $\{e_i\}$, and $V^{-1}$
for the dual space. Both are endowed with the hermitian inner
product with orthonormal basis the vectors $e_0$ and
$(e_i+e_{N-i})/\sqrt{2}$, $i=1,\ldots,N-1$.

Recall the notations of section \ref{PARAM}. For each base surface
$F$ fix an $e$-triangulation $(T,b)$, an ordering of the set
$T^{(2)}$ of $2$-simplices, and let $V(T,b) = \otimes_{t\in
T^{(2)}} V^{\sigma_\pm (t)}$. Given a QHFT bordism
$\Bb=(\tilde{Y},\tilde{L}_\Ff,\rho,\alpha_\pm)$ with a
distinguished flat/charged $\Ii$-triangulation $\Tt(\Bb)$, the
trace tensor in (\ref{deftt}) is a morphism $\Hh_N(\Tt(\Bb)) \in
{\rm Hom}(V(\alpha_-),V(\alpha_+))$, where $V(\alpha_\pm)$ is the
tensor product of isomorphic copies of the spaces $V(T,b)$ over
the (ordered) components of $\alpha_\pm$.
\begin{teo}\label{main} The morphism $\Hh_N(\Bb,h,k)=\Hh_N(\Tt(\Bb))$ does not
depend on the choice of $\Tt(\Bb)$ up to sign and multiplication by
$N$th roots of unity, and $\Hh_N(\Bb,h,k)=\Hh_N(\Bb,h',k')$ if the
mod($N\mz$) reductions of $h$ and $h'$ (resp. $k$ and $k'$) are the
same, that is, if we have $k-k' \in H^1(\tilde{L}_\Ff;\mz)$ and $k-k'=0 \in
H^1(\tilde{L}_\Ff;\mz/N\mz)$, and similarly for $h$ and $h'$. Moreover,
there is no sign ambiguity if we restrict to even valued flattenings
as in Lemma \ref{canonicalflat} (hence with $h=0$). We call
$\Hh_N(\Bb,h,k)$ the {\rm amplitude} of $(\Bb,h,k)$.
\end{teo}
\noindent {\it Proof.} The result up to sign is an immediate
consequence of Proposition \ref{fund_func} and the last claim in
Theorem \ref{teoTOP}. For the dependance with respect to the
mod($N\mz$) reductions of weights, we note that the associated systems
of $N$th roots of moduli (see Remark \ref{univfc}) are connected by
QHG isomorphisms. Indeed, the difference $k-k' \in
H^1(\tilde{L}_\Ff;N\mz)$ coincides with $\gamma(f) -
\gamma(f')/\sqrt{-1}\pi$ and $\gamma(c) - \gamma(c')/\sqrt{-1}\pi$ for
some $f$, $f'$ and $c$, $c'$, and similarly for $h-h'$ and the
$\gamma_2$ maps. By first considering $(k-k')/N$ we can eventually
take the collections of values of $f-f'$ and $c-c'$ in $N\mz$, and
equal. Hence the conclusion follows from (\ref{qvar}) and
(\ref{symqmatdil}). For even valued flattenings, the
claim follows from the fact that Proposition \ref{fund_func} has no
sign ambiguity. For clarity let us state the result with some details.

Consider all the possible branching configurations of $2
\leftrightarrow 3$ QHG transits, up to obvious symmetries. They are
obtained from a single one by composing with the transpositions
$(01)$, $(12)$, $(23)$ and $(34)$ of the vertices (ordered in
accordance with the branching). Any such a transposition changes the
matrix dilogarithm of each tetrahedron by matrix conjugation and
multiplication by a determined scalar factor (see Corollary 5.6 of
\cite{BB2}). Now, there is a prefered ``basic'' $2 \leftrightarrow 3$
QHG transit, for which {\it even} flattenings give no sign ambiguity
in Proposition \ref{fund_func} (see the proof of Theorem 5.7 in
\cite{BB2}). The corresponding branching configuration is defined by
ordering the vertices as $1$ and $3$ on the bottom and top of the
central edge, and $0$, $2$, $4$ counterclockwise as viewed from the
$3$rd vertex. The following table describes for each tetrahedron
$\Delta^i$, opposite to the $i$-th vertex, the scalar factors induced
by the above transpositions (we put $v=\exp(\sqrt{-1}\pi(1-N)/2N)$,
which is $-\zeta^{-(m+1)(1-N)/2}$ in the notations of \cite{BB2}, and
the $c_k^l$ are charge values):
$$\begin{tabular}{||c||c|c||c|c|c||} \hline &
  \rule{0cm}{0.4cm}{$\Delta^1$}& $\Delta^3$ & $\Delta^0$ & $\Delta^2$
  & $\Delta^4$ \\ \hline

\rule{0cm}{0.8cm} \raisebox{0.2cm}{$(01)$} & \raisebox{0.2cm}{$0$} &
\raisebox{0.2cm}{$v^{c_0^3}$} & \raisebox{0.2cm}{$0$} &
\raisebox{0.2cm}{$v^{c_0^2}$} & \raisebox{0.2cm}{$v^{c_0^4}$}\\ \hline

\rule{0cm}{0.8cm} \raisebox{0.2cm}{$(12)$} & \raisebox{0.2cm}{$0$} &
\raisebox{0.2cm}{$v^{c_1^3}$} & \raisebox{0.2cm}{$v^{c_0^0}$} &
\raisebox{0.2cm}{$0$} & \raisebox{0.2cm}{$v^{c_1^4}$} \\ \hline

\rule{0cm}{0.8cm} \raisebox{0.2cm}{$(23)$} &
\raisebox{0.2cm}{$v^{c_1^1}$} & \raisebox{0.2cm}{$0$} &
\raisebox{0.2cm}{$v^{c_1^0}$} & \raisebox{0.2cm}{$0$} &
\raisebox{0.2cm}{$v^{c_0^4}$} \\ \hline

\rule{0cm}{0.8cm} \raisebox{0.2cm}{$(34)$} &
\raisebox{0.2cm}{$v^{c_0^1}$} & \raisebox{0.2cm}{$0$} &
\raisebox{0.2cm}{$v^{c_0^0}$} & \raisebox{0.2cm}{$v^{c_0^2}$} &
\raisebox{0.2cm}{$0$} \\ \hline
\end{tabular}$$
Because of (\ref{preserveq}), we see that the scalars at both sides
are equal. Hence, for any $2 \leftrightarrow 3$ log-branch transit
with even flattenings, the only ambiguity in Proposition
(\ref{fund_func}) is by multiplication by $N$th roots of unity, which
is due to the definition of the function $g$ in (\ref{funcg}). The
conclusion follows as before by using QHG-isomorphisms preserving the
parity of flattenings (whence based on even multiples of the
generators of flat/charge lattices).\hfill$\Box$
\begin{remark}\label{scissors}(Cheeger-Chern-Simons invariants and $\Hh_1$.)
{\rm By the results of \cite{N}, there is an injective
homomorphism from $H_3(PSL(2,\mc);\mz)$ (discrete homology) to a
scissors congruence group $\hat{\mathcal{P}}(\mc)$, such that the
dilogarithm (\ref{symRmatdil}), defined on
$\hat{\mathcal{P}}(\mc)$, restricts to the universal
Cheeger-Chern-Simons class $\hat{C}_2: H_3(PSL(2,\mc);\mz)
\rightarrow \mc/\pi^2\mz$. Hence $\Hh_1$ is a natural extension of
$\exp(\hat{C}_2/i\pi)$, the exponential of a constant times ${\rm
Vol} +iCS$, to classes representing QHFT bordisms. Recently J.
Dupont and C. Zickert produced dilogarithmic formulas for the lift
$\hat{C}_2'$ of $\hat{C}_2$ to $H_3(SL(2,\mc);\mz)$, such that
$\exp(\hat{C}_2'/i\pi)$ coincides with the lift of $\Hh_1$
determined by even flattenings in Theorem \ref{main} \cite{DZ}. We are
indebted to their work for pointing out the existence of such
flattenings.}
\end{remark}
Recall Definition \ref{composition}. Assume that $\Bb'' =
\Bb'*\Bb$ exists, and let $\Tt(\Bb)$ and $\Tt(\Bb')$ be given
weights $(h,k)$ and $(h',k')$, respectively. Then
$\Tt(\Bb')*\Tt(\Bb)$ is a distinguished flat/charged
$\Ii$-triangulation $\Tt(\Bb'')$, with some weight $(h'', k'')$.
(It follows from the Mayer-Vietoris exact sequence
    for $(\Bb'',\Bb',\Bb)$ that even if $h=h'=0$, it can happen
    that $h''\ne 0$. However, if the glued part of the boundary is
    connected, or is a boundary in $\Bb * \Bb'$, then $h=h'=0$ implies
    $h''=0$.)
\begin{prop} \label{functQHFT}{\rm (Functoriality)} For any composition
$\Bb'' = \Bb'*\Bb$ of bordisms, $\Hh_N(\Bb'',h'',k'')$ coincides
with $\Hh_N(\Bb',h',k')\circ \Hh_N(\Bb,h,k)$ up to sign and
multiplication by $N$th roots of unity.
\end{prop}
This is a direct consequence of Theorem \ref{main}. Also, we prove
as in Proposition 4.29 of \cite{BB2}:
\begin{prop}\label{unitarity} {\rm (Polarity)}
Write $\bar{\Bb}$ for the QHFT bordism with opposite orientation and
  complex conjugate holonomy $\bar{\rho}$. Then
  $\Hh_N(\bar{\Bb},-h,-k)$ and $\Hh_N(\Bb,h,k)^*$, the adjoint for the
  hermitian structure of $V(\alpha_\pm)$, coincide up to sign and
  multiplication by $N$th roots of unity.
\end{prop}
In the proof of Theorem \ref{teoTOP} we have seen that the space
$\tilde{W}^+(T,b)$ of ($+$)-Log-$\Ii$-parameters over $W^+(T,b)$
(the disjoint union of spaces $W^+(T,b,\mathcal{C})$) is
isomorphic to $\tilde{\mc}^{-3\chi(F)+2p}$, where $\tilde{\mc}$ is
the universal cover of $\mc \setminus \{0\}$. Similarly, for any
admissible triangulation $\tau$ of a topological cusp with $n$
$2$-simplices, we have the analytic subspace $Def(\tau)$ of
$\widehat{\mc}^n$ made of the $n$-uples of log-branches for the
tetrahedra of $\Ii$-cusps with base triangulation $\tau$, where
$\widehat{\mc}$ is defined in Section \ref{buildbox}. Such
log-branches satisfy the compatibility relations $L_\Tt(e)=0$ at
interior edges.
\begin{defi}\label{defps}{\rm Let $X=(Y,L_\Ff,\alpha_\pm)$ be a marked
topological bordism with $e$-triangulated or admissibly
triangulated boundary components. The {\it phase space} of $X$ is
the (analytic) subset $Def(X)$ in the product of the spaces
$\tilde{W}^+(T,b)$ and $Def(\tau)$ over the components of
$\partial \tilde{Y}$ and $L$, determined by the family of
distinguished flat/charged $\Ii$-triangulations of QHFT bordisms
supported by $(Y,L_\Ff,\alpha_\pm)$.}
\end{defi}
When $Y$ has empty boundary, $Def(X)$ is a generalization of the well
known deformation space of hyperbolic structures supported by ideal
triangulations of $Y\setminus L$, introduced in \cite{Th}, and
recently studied in \cite{Choi} and \cite{Cha}. When $\tilde{Y}$ is
the mapping cylinder of a diffeomorphism $\phi$ of $F$, the amplitudes
of QHFT bordisms supported by $\tilde{Y}$ define a morphism of the
trivial vector bundle
\begin{equation}\label{statebundle}
E(F)^+=E(T,b)^+: \tilde{W}^+(T,b) \times V(T,b) \rightarrow
\tilde{W}^+(T,b). \end{equation}(Note that any two choices of
$e$-triangulations $(T_1,b_1)$ and $(T_2,b_2)$ of $F$ yield
isomorphic bundles $E(T_1,b_1)$ and $E(T_2,b_2)$, with
birationally equivalent bases.) The resulting mapping of sections
of $E(F)$, the {\it states of $F$}, are studied in \cite{B}.
\begin{prop}\label{phase-space}{\rm (Analyticity)} For every $N\geq 1$,
the amplitudes of QHFT bordisms supported by
$X=(Y,L_\Ff,\alpha_\pm)$ vary analytically with the boundary
structure in $Def(X)$, up to sign and multiplication by $N$th
roots of unity.
\end{prop}
This follows immediately from the fact that the matrix
dilogarithms are analytic, together with the fact that any
path in $Def(X)$ lifts to a path of log-branches via the relations
induced by Definition \ref{flat/charge-def} (4).

\subsection{QHFT variants}\label{variant} By varying the
bordism category we can vary the corresponding QFT.
\smallskip

{\bf QHFT$^0$:} Consider the bordism category supported by triples
$(Y,L,\rho)$, where $L$ is an non-empty {\it unframed} tangle in $Y$
and $\rho$ is a $PSL(2,\mc)$-character on the whole of $Y$ (ie. $\rho$
is trivial at the meridians of $L$). In fact, we restrict to
holonomies $\rho$ such that $(Y,L,\rho)$ admits $\Dd$-triangulations
that extend a topological branched ideal triangulation $(T',b')$ of
each boundary component, say $(S,V)$, and for which the link $L$ is
realized as a Hamiltonian subcomplex (hence with no $\Ii$-cusp). In
particular the objects of this bordism category incorporate the
idealization of (necessarily idealizable) cocycles on $(T',b')$, that
represent the restriction of $\rho$ to $S$.  The arguments of Theorem
\ref{main} can be easily adapted to produce tensors $H_N(\Bb,h,k)$
associated to such a bordism $\Bb$, and eventually the so called
QHFT$^0$ variant of quantum hyperbolic field theory.
\medskip

{\bf Fusion of QHFT and QHFT$^0$:} We can consider triples
$(Y,L_\Ff,L^0,\rho)$, where $L=L_\Ff\cup L^0$ is a tangle with a {\it
framed part} $L_\Ff$ and an {\it unframed} one $L^0$.  We also
stipulate that $\rho$ is trivial at each meridian of $L^0$.  For every
object support $(S,V)$, we have a partition $V= V_\Ff \cup V_0$ and we
use ``mixed'' triangulations that looks like an efficient one at
$p\in V_\Ff$ and like an ideal one at $p\in V_0$. A similar mixed
behaviour holds for the adapted $\Dd$-triangulations of such
bordisms. We eventually get tensors still denoted
$\Hh_N(\Bb,h,k)$ giving variants, still denoted QHFT, that extend both
the previous one ($L^0 = \emptyset$), and QHFT$^0$ ($L_\Ff = \emptyset$).
\medskip

{\bf QHFT$^e$:} Let $(Y,L_\Ff,L^0,\rho)$ be as above, and let us
specialize to $\rho$ that, as usual, are trivial at the meridians of
$L^0$, but are {\it not} trivial at the meridians of $L_\Ff$. Now we
use mixed triangulations of each object support $(S,V_\Ff\cup V_0)$
that look like an {\it economic} triangulation (see Remark
\ref{varpar}) at each $p\in V_\Ff$. Concerning the adapted
$\Dd$-triangulations, each component of $L_\Ff$ contributes the
hamiltonian subcomplex with just a copy of the parallel curve
specifying the framing (recall that by using ordinary efficient
triangulations, it contributed with {\it two} parallel curves). We get
tensors now denoted $\Hh^e_N(\Bb,h,k)$, and a variant denoted
QHFT$^e$.
\medskip

Of course, there are no deep structural differences between these
variants; nevertheless each one has its own interest (see also Section
\ref{CLOSED}).

\subsection{Mapping class group representations}\label{MODREP}
Fix $F$. Set $\tilde{Y}_F= F\times [-1,1]$, $\tilde{L}=\partial
F\times [-1,1]$ with trivial vertical framing, and let ${\rm
Mod}(g,r)$ be the mapping class group of $F$, that is, the group of
homotopy classes of orientation-preserving diffeomorphisms of $F$
fixing pointwise each boundary component. Given $\psi_{\pm}:(\pm F,
(T,b,\Ll)) \rightarrow F \times \{\pm 1\}$, put $\psi=
\psi_+^{-1}\psi_-$ and $[\psi]$ for the corresponding element in ${\rm
Mod}(g,r)$. Denote by $\tilde{W}_{[\psi]}$ the {\it mapping torus}
$(F\times [-1,1])/ (x,-1) \sim (\psi(x),1)$ of $\psi$, with tunnel
boundary $\tilde{L}_{[\psi]}$. Let $\rho$ be the conjugacy class of
$PSL(2,\C)$-valued representations of $\pi_1(F\times [-1,1])$
(identified with $\pi_1(F))$) associated to $\Ll$.  To simplify
notations, in all statements of this section we do not mention the
weights (we understand they are fixed).

\begin{lem}\label{CYLtensor} Up to sign and multiplication by
  $N$th roots of unity (denoted "$=_N$") we have :

\smallskip

(1) For any fixed $(T,b,\Ll)$ the amplitudes
$\Hh_N(\tilde{Y},\tilde{L}_F,\rho,(\pm F,(T,b,\Ll)), [\psi_\pm])$
depend only on $[\psi]$. We denote them $\Hh_N([\psi])$.
\smallskip

  (2) $\Hh_N([id])$ is the identity map from $E_N(\alpha_-)$ to
  $E_N(\alpha_+)$, and $\Hh_N([h_2])\circ \Hh_N([h_1]) =_N
  \Hh_N([h_2h_1])$. In particular, for any $[\psi] \in {\rm Mod}(g,r)$
  the QHFT tensor $\Hh_N([\psi])$ is invertible, with inverse
  $\Hh_N([\psi^{-1}])$, and for a homotopically $d$-periodic $[\psi]$
  the QHFT tensor $\Hh_N([\psi])$ is of finite order less than or
  equal to $d$.  \smallskip

  (3) If $\psi(\rho)$ coincides with $\rho$, then ${\rm
    Trace}\bigl(\Hh_N([\psi])\bigr)=_N\Hh_N(\tilde{W}_{[\psi]},
  \tilde{L}_{[\psi]},\rho).$
\end{lem}
\noindent{\it Proof.} Point (1) follows from Theorem \ref{main},
because $$\Hh_N(\tilde{Y},\tilde{L}_F,\rho,(\pm
F,(T,b,\Ll)),[\psi_\pm]) =_N \Hh_N(\tilde{Y},\tilde{L}_F,\rho,(\pm
F,(T,b,\Ll)),[id],[\psi])$$ (the homeomorphism $\psi_-^{-1} \times
{\rm id}$ sends the first mapping cylinder to the second). By
Proposition \ref{functQHFT} we have $\Hh_N^2([id]) =_N
\Hh_N([id])$, so $\Hh_N([id])$ is an idempotent. It is invertible
because the matrix dilogarithms are. Both facts imply the first
claim in (2). The rest is a direct consequence of Proposition
\ref{functQHFT} and formula (\ref{deftt}). \hfill $\Box$
\medskip

The arguments in the proof of Lemma \ref{CYLtensor} imply also
that, letting $[\psi] = [id]$ and $\rho$ fixed, the amplitudes of
any marking variation $(T,b,\Ll) \rightarrow (T_1,b_1,\Ll_1)$ are
invertible. Hence $\Hh_N([\psi])$ is conjugated to the tensor
$\Hh_N(\tilde{Y},\tilde{L}_F,\rho,(\pm F,(T_1,b_1,\Ll_1)),
[\psi_\pm])$. Moreover $\Hh_N(-\tilde{Y},\tilde{L}_F,\rho,(\mp
F,(T,b,\Ll)), [\psi_\pm])$, the amplitude with reversed
orientation, clearly coincides with $\Hh_N([\psi]^{-1})$. Using
Proposition \ref{unitarity} we deduce:
\begin{cor}\label{mcgrep} For any fixed $\rho\in \Rr(g,r)$, the homomorphisms
$\psi \mapsto \Hh_N([\psi])$ induce a conjugacy class of
  linear representations of ${\rm Mod}(g,r)$, well-defined up to sign and
  multiplication by $N$th roots of unity. For $SL(2,\mr)$-valued
  characters $\rho$ these representations are unitary.
\end{cor}

\subsection{Tunneling the ($+$)/($-$) states}\label{+-versus} We use
($+$)-Log-$\Ii$-parameters to define the QHFT because of the existence
of strongly idealizable cocycles on QHFT bordism triangulations, which
makes functoriality easy to check. Here we exhibit a family of tensors
correlating the ($\pm$)-Log-$\Ii$-parameters, thus recovering, in
particular, the direct and nice interpretation of boundary structures
having non trivial holonomy at the punctures in terms of pleated
hyperbolic surfaces (see Section \ref{Ipar}). These tensors are also
used in Section \ref{FIBERED}.
\smallskip

For any base surface $F$ with an $e$-triangulation $(T,b)$, let $Z(F)$
be the pseudo-manifold obtained by collapsing to a point each boundary
annulus of the cylinder $C(F)=F\times [-1,1]$.  Recall the bundle
$E(F)^{+}$ in (\ref{statebundle}), and consider similarly
$E(F)^{-}:\tilde{W}^-(T,b) \times V(T,b) \rightarrow
\tilde{W}^-(T,b)$. We have:
\begin{prop}\label{tunnel+-} There exists a canonical family $\mathcal{F}$
of flat/charged $\Ii$-triangulations covering a portion of $Z(F)$,
with invertible trace tensor $\Hh_N(\mathcal{F}): E(F)^{+}
\rightarrow E(F)^-$.
\end{prop}
\noindent {\it Proof.} Orient $C(F)$ so that $\pm F$ is identified
with $F\times \{\pm 1\}$. Let $P(T,b)$ be the cell decomposition
of $C(F)$ made by the prisms with base the $2$-simplices of $T$.
Orient all the ``vertical'' (ie. parallel to $[-1,1]$) edges of
$P(T,b)$ towards $+F$. For every abstract prism $P$, every
vertical boundary quadrilateral $R$ has both the two horizontal
and the two vertical edges endowed with parallel orientations. So
exactly one vertex of $R$ is a source (that belongs to $-F$), and
exactly one is a pit (that belongs to $+F$). Triangulate each $R$
by the oriented diagonal going from the source to the pit. Finally
extend the resulting triangulation of $\partial P$ to a
triangulation of $P$ made of $3$ tetrahedra, by taking the cone
from the $b$-first vertex of the bottom base triangle of $P$ (note
that no further vertices nor further edges have been introduced).
Repeating this for every prism, we get a branched distinguished
triangulation $(C(T,b),H)$ of $C(F)$, where the vertical edges
make the Hamiltonian tangle $H$. As in the proof of Theorem
\ref{teoTOP} there exists integral charges on $(C(T,b),H)$.

Let $F\times [-1,3]$ be triangulated by two adjacent copies of
$C(T,b)$, glued each to the other at $F\times \{+1\}$. For any $z \in
\Zz_I(T,b)^+$, consider the unique cocycle $C_0(z)$ on the composition
$C(T,b)*C(T,b)$ that: extends $z\cup \Psi(z)$, given on $(F\times
\{-1\})\cup (F\times \{3\})$; takes the value $P$ of (\ref{convert})
on each vertical edge contained in $F\times [1,2]$; takes the value
$1$ on each vertical edge contained in $F\times [-1,1]$. Perturb
$C_0(z)$ with a $0$-cochain $s$ that: takes the value $1$ on $(F\times
\{-1\})\cup (F\times \{3\})$; takes values in $PB^+(2,\C)$ at each
vertical boundary annulus; restricts to an idealizable cocycle on
$F\times [-1,1]$ and to a {\it maximally idealizable} cocycle on
$F\times [1,3]$ (see Definition \ref{triideal}). Finally, glue
$\Ii$-cusps to the idealization.  Note that the only non idealizable
tetrahedra are those in the star of a boundary edge of the
triangulation of $F\times \{3\}$.  Lemma \ref{canonicalflat} gives
flattenings for the idealizable tetrahedra.

Look at the ideal triangulation $(T',b')$ of $S\setminus V$
corresponding to the copy of the triangulation $(T,b)$ for the
boundary component $F\times \{3\}$. For every cochain $s$ as above and
every edge $e$ of $T'$ we have two complex numbers: the
($-$)-exponential $\Ii$-parameter $W^-(\Psi(z))(e)$, and, as in
(\ref{totw}), the total product $\Ww(e)$ of the cross-ratio moduli at
the edges of $T$ that enter the definition of
$W^-(\Psi(z))(e)$. Recall that there are two distinct such edges only when $e$
contributes to make a marked corner. It is possible to normalize $s$
so that for every edge $e$ of $T'$ we have $W^-(\Psi(z))(e) =
\Ww(e)^{-1}$. 

Varying the cocycle $z \in \Zz_I(T,b)^+$, this choice determines the
family $\mathcal{F}$ in the statement. By perturbing the initial
cocycle $z$ with $0$-cochains $t$ with values in $PSL(2,\mc)\setminus
PB^+(2,\mc)$, the same construction leads to families $\mathcal{F}_t$
of flat/charged $\Ii$-triangulations covering the whole of $F \times
[-1,3]$. 

Note that for suitable flat/charges ($f_1=0$ in (\ref{symRmatdil}) and
$c_1=0$ in (\ref{symqmatdil})) the matrix dilogarithms have
well-defined finite limits when the cross-ratio modulus $w_0
\rightarrow 0$. From the symmetry relations of the matrix dilogarithms
(see \cite{BB2}, Corollary 5.6), this is true more in general for any
degenerating sequence of $\Ii$-tetrahedra, that is when $w_0$ goes to
$0$, $1$ or $\infty$. Now, we can choose in a continuous way the
flat/charges of $\mathcal{F}_t$ so that they satisfy the above
constraints on the tetrahedra of $\mathcal{F}_t$ that become non
idealizable in $\mathcal{F}$, when $t \rightarrow {\rm id}$. Then
$\textstyle \Hh_N(\mathcal{F}):=\lim_t \Hh_N(\mathcal{F}_t)$
exists. As in Lemma \ref{CYLtensor} (2) we see that
$\Hh_N(\mathcal{F}_t)$ is invertible, with inverse
$\Hh_N(-\mathcal{F}_t)$. Since $\textstyle \Hh_N(\mathcal{F}) \circ
\Hh_N(-\mathcal{F})= \lim_t (\Hh_N(\mathcal{F}_t)\circ
\Hh_N(-\mathcal{F}_t))=_N{\rm id}$ (Proposition \ref{phase-space}), we
deduce that $\Hh_N(\mathcal{F})$ is invertible.  \hfill$\Box$
\begin{figure}[ht]
\begin{center}
\includegraphics[width=5cm]{Tcyl.eps}
\caption{\label{tcusp} Pasting opposite vertical sides yields an
instance of $C(T,b)$ for the once-punctured torus $S$, based on an
economic triangulation of $S$ as in Remark \ref{varpar}.}
\end{center}
\end{figure}
